%% ============================================================
%% Chapter 8 — Discussion
%% ============================================================
\chapter{Discussion}
\label{chap:discussion}

\section{Information Value of Retrieval Confidence vs Reader Confidence}
\label{sec:disc_info_value}

A central question posed at the outset of this investigation was whether sparse retrieval signals alone provide sufficient discriminative power to guide abstention decisions in retrieval-augmented question answering. Our empirical results provide a decisive answer: \textbf{BM25 scores alone are fundamentally inadequate for gating answerability.}

As shown by the comparison between the \emph{Always Answer} baseline and the \emph{BM25 Threshold} strategy in Table~\ref{tab:strategy_comparison}, thresholding on the top BM25 score yielded an answerability $F_1$ of $0.6486$, representing virtually no improvement over the naïve baseline of $0.6505$. More strikingly, the optimal threshold tuned on the validation set resulted in abstaining on only a single question in the entire test set (achieving $0.998$ coverage and a false-answer rate of $1.0000$). 

The underlying reason for this failure stems from the nature of term-matching retrieval. BM25 scores reflect surface-level lexical overlap and inverse document frequency statistics, not factual sufficiency or semantic entitlement. SQuAD~2.0 unanswerable questions were specifically engineered by human crowdworkers to exhibit high lexical overlap with their paired passages while subtly altering key entities, dates, or relational constraints. Consequently, unanswerable queries often trigger retrieval scores that are indistinguishable from, or even higher than, answerable queries (as illustrated in Figure~\ref{fig:bm25_score_distribution}).

In sharp contrast, the multi-feature \emph{Answerability Classifier} achieved an $F_1$ score of $0.8543$, reducing the false-answer rate from $1.0000$ to $0.1776$. This leap in performance confirms that combining retrieval-side statistics with reader-side confidence metrics (\texttt{qa\_score}, \texttt{qa\_no\_answer\_probability}, and predicted answer length) captures critical orthogonal information. While BM25 signals whether topically related text has been surfaced, the reader component models whether a coherent, well-supported entity span answering the precise relation in the query actually exists within that text.

\section{The Precision-Recall and Risk-Coverage Trade-off}
\label{sec:disc_tradeoff}

In selective prediction frameworks, there is an unavoidable trade-off between coverage (the proportion of queries on which the system agrees to answer) and precision (the accuracy or factual validity of the answers provided).

Figure~\ref{fig:pr_curve} and the strategy comparison highlight this dynamics clearly:
\begin{itemize}
  \item Under the \emph{Always Answer} paradigm, coverage is $100\%$, but the risk of hallucination or false assertions is maximised. In our balanced test set, $50\%$ of answers were presented on unanswerable questions, leading to a false-answer rate of $100\%$ and strict Exact Match of only $1\%$.
  \item By tuning the decision threshold to $\delta^* = 0.42$ on validation data, the Answerability Classifier operates at an operating point retaining $85.0\%$ coverage and $84.0\%$ recall, cutting false abstentions to $16.0\%$ while identifying queries where reader confidence collapses.
  \item In contrast, the BM25 threshold heuristic ($\tau_{\text{BM25}} = 13.15$) reduces false answers to $72.0\%$ and achieves the highest overall strict Exact Match ($15\%$), but incurs a heavy false-abstention rate of $30.0\%$, discarding viable answers.
\end{itemize}

In mission-critical domains (such as medical diagnosis, legal statutory search, or automated technical compliance), producing an incorrect answer carries severe legal and safety consequences, whereas an explicit refusal (``insufficient information available'') allows a human operator to intervene. In such contexts, operating at an even more conservative threshold (e.g., $\delta \ge 0.50$) would further depress false answers at the cost of lower coverage.

\section{Methodological and Dataset Limitations}
\label{sec:disc_limitations}

While the experimental findings are robust, several methodological caveats must be acknowledged:

\begin{enumerate}
  \item \textbf{SQuAD Passage-Level Answerability Proxy Assumption:} The ground-truth label \texttt{is\_answerable} in SQuAD~2.0 is defined strictly relative to the original paired Wikipedia paragraph. In a retrieval-augmented setting, BM25 searches across the broader partition corpus. SQuAD 2.0 unanswerable queries were deliberately authored by crowdworkers to share dense lexical overlap with candidate passages, making sparse BM25 scoring particularly vulnerable to adversarial distractors. Furthermore, passage-level answerability labels do not directly tell us if an alternative retrieved passage contains an accidental answer.
  \item \textbf{Experimental Scale (500 Total Examples):} To maintain computational efficiency, leak-free article partitioning, and complete reproducibility on standard CPU hardware, the complete QA pipeline was trained and evaluated on a rigorously balanced 500-example study (300 training, 100 validation, 100 test). The partition preserves 50\% class balance and strict article separation (seed 42), ensuring statistical rigor without inflating claims to the full multi-thousand SQuAD corpus.
  \item \textbf{Extractive Reader Constraint:} The experimental setup utilizes an extractive model (\texttt{deepset/roberta-base-squad2}) rather than a generative autoregressive LLM (e.g., LLaMA-3 or Mistral). While extractive readers provide strictly bounded spans and well-calibrated start/end logits, modern RAG systems frequently rely on generative synthesis. Evaluating abstention mechanisms across generative backbones remains an essential follow-up.
\end{enumerate}

\section{Implications for Real-World RAG Deployment}
\label{sec:disc_deployment}

Our findings offer direct architectural recommendations for production RAG pipelines:
\begin{itemize}
  \item \textbf{Never gate on vector or BM25 similarity alone:} In practice, similarity thresholds (e.g., cosine similarity $> 0.75$) are frequently used as heuristics to reject queries. Our work empirically proves that retrieval similarity does not imply factual sufficiency.
  \item \textbf{Deploy a dedicated gating verifier:} An inexpensive logistic regression or lightweight MLP that ingests cross-encoder reranker scores, reader span logits, and lexical overlap features provides a reliable, sub-millisecond abstention layer before incurring expensive generation or exposing users to hallucinations.
\end{itemize}
